\documentclass[11pt,a4paper]{ctexart}
\usepackage[margin=2.2cm]{geometry}
\usepackage{amsmath,amssymb,amsthm,bm,mathtools}
\usepackage{booktabs,longtable,array,multirow}
\usepackage{enumitem}
\usepackage{hyperref}
\usepackage[dvipsnames]{xcolor}
\usepackage[most]{tcolorbox}
\usepackage{graphicx}
\usepackage{float}
\usepackage{titlesec}
\usepackage{fancyhdr}
\usepackage{lastpage}
\usepackage{setspace}
\usepackage{caption}
\usepackage{tabularx}
\usepackage{makecell}
\hypersetup{colorlinks=true,linkcolor=MidnightBlue,urlcolor=MidnightBlue,citecolor=MidnightBlue}
\setlength{\parindent}{2em}
\setlength{\parskip}{0.3em}
\onehalfspacing
\pagestyle{fancy}
\setlength{\headheight}{15pt}
\fancyhf{}
\fancyhead[L]{SABER-PLR 设计文档}
\fancyhead[R]{\thepage/\pageref{LastPage}}
\fancyfoot[C]{}
\titleformat{\section}{\Large\bfseries\color{MidnightBlue}}{\thesection}{0.6em}{}
\titleformat{\subsection}{\large\bfseries\color{NavyBlue}}{\thesubsection}{0.6em}{}
\titleformat{\subsubsection}{\normalsize\bfseries\color{RoyalBlue}}{\thesubsubsection}{0.6em}{}
\newcommand{\E}{\mathbb{E}}
\newcommand{\I}{\mathbf{1}}
\newcommand{\R}{\mathbb{R}}
\newcommand{\softmax}{\operatorname{softmax}}
\newcommand{\clip}{\operatorname{clip}}
\newcommand{\Var}{\operatorname{Var}}
\newcommand{\sigmoid}{\operatorname{sigmoid}}
\newcommand{\norm}[1]{\left\lVert #1 \right\rVert}
\newcommand{\argmax}{\operatorname*{argmax}}
\newcommand{\argmin}{\operatorname*{argmin}}
\captionsetup{font=small,labelfont=bf}

\begin{document}

\begin{titlepage}
\centering
\vspace*{2.0cm}
{\Huge\bfseries SABER-PLR：面向 hardest OOD 的\\[0.4em]严重度与分支感知高效课程重放框架\par}
\vspace{1.0cm}
{\Large Severity- and Branch-aware Efficient Replay Prioritization over Levels\par}
\vspace{1.6cm}
\begin{tcolorbox}[colback=blue!3,colframe=MidnightBlue,boxrule=0.8pt,width=0.92\textwidth]
\textbf{文档目的}\quad 给出一套\textbf{按工程现实逐步实现}的新算法家族，用于在你们现有 \texttt{PPO\_NEW-v3/v3.1} 主干之上，重构当前 PLR\_UED 外层流程。设计目标有三点：
\begin{enumerate}[leftmargin=2.2em,itemsep=0.25em]
    \item 保留现有 PLR/UED 在 \(O\_\{10,90\}\) hardest OOD 上对高严重度边界切片的有效修复能力；
    \item 减少外层昂贵交互次数，并在完成 runtime profiling 后争取显著压低总耗时，而不是先验承诺固定的 wall-clock 降幅；
    \item 在 telemetry 与特征导出链路打通的前提下，逐步让外层分布调度感知 \texttt{v3} 的三路结构 \((o,\,meta,\,\Delta o)\)；若链路未打通，则 branch-aware 只能作为后续增强，而不是第一批必做项。
\end{enumerate}
\end{tcolorbox}
\vfill
{\large 版本：v1.1\par}
{\large 生成日期：\today\par}
\end{titlepage}

\tableofcontents
\newpage

\section{问题重述与设计出发点}

\subsection{当前最关键的经验事实}
根据你们当前\textbf{已完成且可核对}的基线与 trace-level 切片分析，现阶段可以较稳地说：
\begin{enumerate}[leftmargin=2.2em,itemsep=0.35em]
    \item 在 strongest OOD 分布 \(O\_{10,90}\) 上，普通 PPO 的平均回报长期停留在约 \(0.5\) 的低平台，这至少说明当前 hardest OOD 下存在稳定失效模式，而不只是单次随机波动；
    \item \texttt{PPO\_NEW-v3/v3.1} 作为 inner backbone 在 \(F1\) 与部分 \(G\) 分布上有收益，但对 \(O\_{10,90}\) 并未稳定超过普通 PPO；
    \item 当前 outer 系方法里，主结果上\textbf{明确跑赢 PPO 的只有 PLR\_UED}；而且它的提升不是全局均匀抬升，而是集中修复了 \textbf{高严重度 \(severity=5/6\) \(\times\) removal} 的边界决策，并伴随少量 insertion 侧收益；
    \item 因此当前证据更支持如下\textbf{优先级判断}，而不是绝对定理：
     \[
     \text{训练分布选择/调度效应} > \text{inner 表征归纳偏置效应} > \text{outer 控制器复杂度}.
     \]
\end{enumerate}

\subsection{当前必须承认的工程与证据边界}
为了避免后文把“研究设想”误写成“现成能力”，这里先把边界说清楚：
\begin{enumerate}[leftmargin=2.2em,itemsep=0.35em]
    \item \(O\_{10,90}\) 上关于 hardest slice 的机制证据最强，但目前仍以有限 seeds 与有限完整 run 为主，因此适合指导算法设计，不宜写成过度确定的理论结论；
    \item 当前最可信的 trace-level 信号是 \(severity\)、stream、语义动作、reward，以及 implement 路径上部分可用的 \(p(a=1)\)；训练路径并未稳定导出完整 \(p(a=1)\)；
    \item 当前 \texttt{PPO\_NEW-v3} 代码主干稳定暴露的是融合后的最终特征 \(h\)，并未稳定导出三路 branch tensor 或 branch 级梯度，因此 branch-aware 设计目前只能视为\textbf{可选研究扩展}；
    \item 当前 outer action space 本质上是离散 level 网格，而不是连续可微参数空间；因此任何“邻域编辑”都必须遵守当前 generator 真正支持的离散候选，而不是假设连续局部扰动天然存在。
\end{enumerate}

\subsection{我们真正要修复的坏点}
记每个 implement 样本为
\[
\tau_t=(s_t, a_t, r_t, sev_t, str_t),
\]
其中 \(sev_t\) 表示严重度，\(str_t\in\{\texttt{removal},\texttt{insertion}\}\) 表示当前 stream。由已有切片结果可知，hardest OOD 的核心失败模式并非“所有样本都差”，而是集中在集合
\[
\mathcal{H}=\{t\mid sev_t\ge \tau_H,\; str_t=\texttt{removal}\},
\]
其中默认推荐 \(\tau_H=5\)。

普通 PPO 在该集合上几乎完全塌缩到保守动作 \(a=0\)；而 PLR\_UED 的收益则主要来自：
\begin{enumerate}[leftmargin=2.2em,itemsep=0.25em]
    \item 在 \(\mathcal{H}\) 上恢复少量但高质量的 \(a=1\)；
    \item 在 \(sev=5\) 左右保留更好的“继续保守也要保守得准确”的能力；
    \item 不追求全局激进，而追求 \textbf{selective action1 recovery}。
\end{enumerate}

因此，新算法不应再是“更通用的 outer RL”，而应是：
\begin{quote}
\textbf{围绕 hardest slice 进行定点课程调度的轻量级课程重放器。}
\end{quote}

\subsection{新算法总名称与版本路线}
本文给出一个从简单到复杂、便于逐版实现的算法家族，统称
\begin{center}
\fbox{\parbox{0.9\textwidth}{\centering\textbf{SABER-PLR: Severity- and Branch-aware\ Efficient Replay Prioritization over Levels}}}
\end{center}
并拆为五个递进版本：
\begin{center}
\begin{tabular}{p{2.1cm}p{3.0cm}p{8.6cm}}
\toprule
版本 & 名称 & 核心新增点 \\
\midrule
V0 & PLR-Lite & 最小可替代版：保留 replay 与 learnability，去掉重 outer RL，仅做 score-based level selector。 \\
V1 & SR-PLR & 加入 \textbf{severity-aware + removal-aware} 定向打分，主攻 hardest slice。 \\
V2 & SBR-PLR & \textbf{可选研究扩展}：先做 V2a 的 proxy-only branch-aware，再视 telemetry 情况升级到 V2b。 \\
V3 & BSR-PLR & 在 V1 之上加入 \textbf{learnability 硬门 + 两阶段预算控制} 与 cost-aware utility；这是最有希望先落地降时的版本。 \\
V4 & SABER-PLR & 完整增强版：在 V1/V3 稳定后，再考虑邻域编辑课程、adaptive new/replay mixing、轻量 selector 等机制。 \\
\bottomrule
\end{tabular}
\end{center}

\section{记号、主干网络与总体设计原则}

\subsection{内层 backbone：保持 PPO\_NEW-v3/v3.1 稳定主干不动}
我们不修改当前 PPO 主优化器，只把 \texttt{v3/v3.1} 作为稳定 backbone。给定时刻 \(t\) 的输入，定义
\[
 x_t=[o_t,\,b_t,\,a_{t-1},\,r_{t-1},\,\Delta o_t],
\]
其中：
\begin{itemize}[leftmargin=2.2em,itemsep=0.2em]
    \item \(o_t\)：当前局部观测；
    \item \(b_t\)：阶段位或 meta 状态；
    \item \(a_{t-1},r_{t-1}\)：最短闭环行为反馈；
    \item \(\Delta o_t\)：短时变化量。
\end{itemize}

在 \texttt{v3} 中，这些信号被拆成三路：
\[
 z_t^{(o)}=f_o(o_t),\quad z_t^{(m)}=f_m(b_t,a_{t-1},r_{t-1}),\quad z_t^{(\Delta)}=f_{\Delta}(\Delta o_t),
\]
然后拼接或融合成
\[
 z_t = \phi\big(z_t^{(o)}, z_t^{(m)}, z_t^{(\Delta)}\big),
\]
再经过 TemporalConv 或轻量时序聚合器得到最终策略特征 \(h_t\)，最后进入 actor/critic 头。

本设计中，\textbf{inner PPO 训练逻辑不改}，新算法只改 outer 的分布选择、候选重放、评估预算与评分信号。

\subsection{外层 level 定义}
一个外层 level 记为
\[
 g=(\mu_A,\mu_B,p,n,pattern),
\]
分别表示：
\begin{itemize}[leftmargin=2.2em,itemsep=0.2em]
    \item \(\mu_A,\mu_B\)：两类关键事件或上下文参数的均值；
    \item \(p\)：混合比例或切换比例；
    \item \(n\)：当前 generator 所支持的离散文件规模或离散长度档位；
    \item \(pattern\)：当前 generator 所支持的离散模式。
\end{itemize}
在当前项目里，outer action space 更接近\textbf{离散 level 网格}，而不是连续可微参数控制。第一版实现应优先复用现有离散 key、现有 replay 机制与现有候选选择器，再讨论局部邻域编辑。

\subsection{四条总原则}
\begin{enumerate}[leftmargin=2.2em,itemsep=0.25em]
    \item \textbf{定点修 hardest slice。} 重点不是全局均值，而是 \(severity\ge 5\) 的 removal 边界。 
    \item \textbf{选择性恢复 action1。} 目标不是全局提高 \(action1\_rate\)，而是只在“该出手”的高严重度边界上提高高置信度 \(a=1\)。
    \item \textbf{removal/insertion 非对称。} 同一个二值动作在两个 stream 里的语义相反，因此 outer score 必须拆流统计。
    \item \textbf{少算昂贵样本。} 你们的瓶颈主要是 ALNS 相关外层交互成本，因此必须用两阶段预算与候选筛选减少 full evaluation 次数。
\end{enumerate}

\section{统一评分框架：从 PLR 到 SABER-PLR}

\subsection{基础 PLR/UED 项}
对任意候选 level \(g\)，记冻结策略或当前参考策略在该 level 上的表现为 \(J_{frozen}(g)\)，记短训或最新统计得到的表现为 \(J(g)\)。

沿用 PLR/UED 的核心思想，首先定义一个 \textbf{可学习门}：
\begin{equation}
L(g)=\exp\left(-\frac{(J_{frozen}(g)-c)^2}{2\sigma_L^2}\right).
\label{eq:learnability}
\end{equation}
当 \(J_{frozen}(g)\) 太高时，说明太简单；太低时，说明太难。\(L(g)\) 迫使采样集中到“难但可学”的边界附近。

再定义学习进度项：
\begin{equation}
LP(g)=\big|J(g)-\bar J_{ref}(g)\big|,
\label{eq:lp}
\end{equation}
其中 \(\bar J_{ref}(g)\) 可以是上一轮、EMA 基线或历史缓冲统计。由于当前任务奖励近似二值，绝对回报对细粒度能力变化不够敏感，因此 \(LP(g)\) 比单纯 \(1-J\) 更重要。

最后，沿用 replay buffer 的优先级更新：
\begin{equation}
q_g \leftarrow (1-\beta_q)q_g + \beta_q\,score_g,
\label{eq:q_update}
\end{equation}
采样概率为
\begin{equation}
P(g)=\frac{(q_g+\varepsilon_q)^{\alpha_q}}{\sum_{g'}(q_{g'}+\varepsilon_q)^{\alpha_q}}.
\label{eq:priority_sampling}
\end{equation}

\subsection{hardest-slice 定向项}
定义 hardest removal slice 与 hardest insertion slice 为
\begin{equation}
\mathcal{H}_{rem}(g)=\{t\mid sev_t\ge \tau_H,\; str_t=\texttt{removal}\},
\label{eq:hard_slice}
\end{equation}
\begin{equation}
\mathcal{H}_{ins}(g)=\{t\mid sev_t\ge \tau_H,\; str_t=\texttt{insertion}\}.
\label{eq:hard_slice_ins}
\end{equation}
默认 \(\tau_H=5\)。为简化记号，后文若未特别说明，\(\mathcal{H}(g)\equiv \mathcal{H}_{rem}(g)\)，即\textbf{主 score 默认只指向 removal hardest slice}。同时定义 easy slice 为
\begin{equation}
\mathcal{E}(g)=\{t\mid sev_t\le \tau_E\},\qquad \tau_E=3.
\label{eq:easy_slice}
\end{equation}

当前证据表明，PLR\_UED 的有效增益并不只是“多做 action1”，还包括 hardest slice 上继续保守时的正确性。因此至少需要五个定向统计量。

\paragraph{(1) hardest-slice 总体正确性}
\begin{equation}
Q_{hard}(g)=\E_{t\in \mathcal{H}(g)}\big[r_t\big].
\label{eq:Qhard}
\end{equation}
它反映的是 hardest removal slice 的\textbf{总体可处理程度}。这一项必须保留，因为现有 trace 证据说明 PLR\_UED 的收益不只来自恢复少量 \(a=1\)，也来自高严重度下更准确的保守决策。

\paragraph{(2) 高严重度移除边界恢复分数}
\begin{equation}
R_{hard}(g)=\E_{t\in \mathcal{H}(g)}\big[\I(a_t=1)\,r_t\big].
\label{eq:Rhard}
\end{equation}
它度量：在 hardest slice 上，是否恢复了真正有用的 \(a=1\)。但若单独使用这一项，容易把 score 推向“过度追求 action1”，因此必须与 \(Q_{hard}\) 联合使用。

\paragraph{(3) 选择性置信分数}
记策略在 removal 上预测 \(a=1\) 的概率为 \(p_t=\pi_\theta(a=1\mid h_t)\)。定义
\begin{equation}
C_{sel}(g)=\E_{t\in \mathcal{H}(g)}\big[p_t\,\I(a_t=1)\,r_t\big]-\lambda_{fp}\E_{t\in \mathcal{E}(g)}\big[\I(a_t=1)\big].
\label{eq:csel}
\end{equation}
第一项奖励“高严重度下高置信度的正确出手”，第二项惩罚 easy slice 上的无谓激进。

若当前 probe 链路\textbf{无法稳定导出} \(p_t\) 或 \(p(a=1)\)，则应退化为
\begin{equation}
\tilde C_{sel}(g)=\E_{t\in \mathcal{H}(g)}\big[\I(a_t=1)\,r_t\big]-\lambda_{fp}\E_{t\in \mathcal{E}(g)}\big[\I(a_t=1)\big],
\label{eq:csel_fallback}
\end{equation}
或直接令 \(w_c=0\)。这里不能把“理论上可算”误当成“当前工程上已稳定可算”。

\paragraph{(4) easy-slice 保护项}
\begin{equation}
P_{easy}(g)=\E_{t\in \mathcal{E}(g)}\big[\I(a_t=0)\,r_t\big].
\label{eq:peasy}
\end{equation}
它用于抑制为了追求 hardest slice 而伤害 easy slice 的副作用。

\paragraph{(5) insertion 弱辅助监控项}
由于 insertion 与 removal 语义相反，而且当前 trace 证据显示主收益集中在 removal hardest slice，因此 insertion 在 V1/V3 中\textbf{不与 removal 对称建模}。定义一个弱辅助监控量
\begin{equation}
M_{ins}(g)=\E_{t\in \mathcal{H}_{ins}(g)}\big[r_t\big].
\label{eq:mins}
\end{equation}
它只表示：若某些 candidate level 顺带改善了高严重度 insertion slice，可给予\textbf{很小}的辅助奖励。但第一版完全可以取 \(w_i=0\)，把 \(M_{ins}\) 仅作为监控指标，而不进入主优化目标。这样可以明确避免把 insertion 与 removal 完全对称处理。

\subsection{可选研究扩展：branch-aware learnability}
这一节描述的是\textbf{后续可选扩展}，不是第一批实现的前置条件。当前代码与日志链路尚未稳定导出三路 branch tensor 与 branch 级梯度，因此这里更适合作为中后期研究方向，而不是首版工程规格。其目标不是强迫三路绝对均衡，而是：
\begin{quote}
\textbf{避免长期只有一条分支承担 hardest slice 决策，而另外两条分支几乎失活或不可学习。}
\end{quote}

\paragraph{V2a：proxy-only branch-aware（当前更现实的首版）}
V2a 只依赖 forward-time 的 activation-level 或 output-level 信号，不要求 per-branch gradient telemetry。记 \(u_t^{(b)}=\norm{z_t^{(b)}}_2\)，定义分支活性占比
\begin{equation}
 a_{t,\mathrm{act}}^{(b)} = \frac{u_t^{(b)}}{\sum_{b'} u_t^{(b')}+\varepsilon_a},
 \qquad b\in\{o,m,\Delta\}.
 \label{eq:branch_act_ratio}
\end{equation}
再定义一个轻量 ablation sensitivity。记 \(\ell_t^{(-b)}\) 为将第 \(b\) 路分支置零、掩蔽或替换为固定基线后得到的当前动作 logit，则
\begin{equation}
 \delta_t^{(b)} = \big|\ell_t-\ell_t^{(-b)}\big|,\qquad
 \tilde\delta_t^{(b)} = \frac{\delta_t^{(b)}}{\sum_{b'}\delta_t^{(b')}+\varepsilon_\delta}.
 \label{eq:branch_ablation}
\end{equation}
在 hardest slice 上求均值得到
\begin{equation}
 A_{b,\mathrm{act}}(g)=\E_{t\in\mathcal{H}(g)}\big[a_{t,\mathrm{act}}^{(b)}\big],\qquad
 S_{b,\mathrm{abl}}(g)=\E_{t\in\mathcal{H}(g)}\big[\tilde\delta_t^{(b)}\big].
 \label{eq:branch_proxy_stats}
\end{equation}
于是 V2a 的 proxy 分数定义为
\begin{equation}
S_{branch}^{V2a}(g)=\sum_{b\in\{o,m,\Delta\}}\log\big(\varepsilon_b + A_{b,\mathrm{act}}(g)S_{b,\mathrm{abl}}(g)\big)
- \eta_b\,\Var\big(A_{o,\mathrm{act}}S_{o,\mathrm{abl}}, A_{m,\mathrm{act}}S_{m,\mathrm{abl}}, A_{\Delta,\mathrm{act}}S_{\Delta,\mathrm{abl}}\big).
\label{eq:sbranch_v2a}
\end{equation}
这一路线只需要\textbf{前向可见的 branch output 与简易遮挡实验}，因此是当前更现实的第一实现版本。

\paragraph{V2b：gradient-aware branch-aware（telemetry 打通后的增强版）}
若后续 telemetry 打通，记 \(\ell_t\) 为策略对当前动作的 logit。对每个 hardest slice 样本，定义三路分支的 \textbf{相对敏感度}：
\begin{equation}
 s_t^{(b)} = \norm{\frac{\partial \ell_t}{\partial z_t^{(b)}}}_2\cdot \norm{z_t^{(b)}}_2,
 \qquad b\in\{o,m,\Delta\}.
 \label{eq:sensitivity}
\end{equation}
归一化得到三路相对贡献：
\begin{equation}
 a_t^{(b)} = \frac{s_t^{(b)}}{\sum_{b'} s_t^{(b')}+\varepsilon_a}.
 \label{eq:branch_ratio}
\end{equation}
在 hardest slice 上求均值得到
\begin{equation}
 A_b(g)=\E_{t\in\mathcal{H}(g)}\big[a_t^{(b)}\big].
 \label{eq:Ab}
\end{equation}

仅有活性还不够，还需要“可学习”。因此再定义 probe 期间每条分支的梯度可学习度：
\begin{equation}
 G_b(g)=\E_{t\in\mathcal{H}(g)}\left[\norm{\nabla_{\theta_b}\mathcal{L}^{probe}_{ppo,t}}_2\right],
 \label{eq:Gb}
\end{equation}
并归一化为 \(\tilde G_b(g)=G_b(g)/(\sum_{b'}G_{b'}(g)+\varepsilon_G)\)。

最后构造更强的 gradient-aware 分数：
\begin{equation}
S_{branch}^{V2b}(g)=\sum_{b\in\{o,m,\Delta\}}\log\big(\varepsilon_b + A_b(g)\tilde G_b(g)\big) - \eta_b\,\Var\big(A_o\tilde G_o, A_m\tilde G_m, A_\Delta\tilde G_\Delta\big).
\label{eq:sbranch_v2b}
\end{equation}
其含义是：
\begin{itemize}[leftmargin=2.2em,itemsep=0.2em]
    \item 每条分支都要在 hardest slice 上真正参与，且存在可学习梯度；
    \item 若只有一路过强、两路失活，则方差项给出惩罚；
    \item 若某条分支在当前候选 level 上几乎不起作用，则该 level 不应被过度优先。
\end{itemize}

在总分中统一记
\begin{equation}
S_{branch}(g)=
\begin{cases}
S_{branch}^{V2a}(g), & \text{若当前采用 V2a/proxy-only 模式},\\
S_{branch}^{V2b}(g), & \text{若 telemetry 已打通且采用 V2b 模式}.
\end{cases}
\label{eq:sbranch}
\end{equation}

\paragraph{实现前提与保守落地方式}
若要真正落地 branch-aware，建议按两级推进：
\begin{enumerate}[leftmargin=2.2em,itemsep=0.2em]
    \item \textbf{V2a 先行。} 只要求 backbone 在 forward 阶段能导出 \(z^{(o)},z^{(m)},z^{(\Delta)}\) 或其等价 proxy，并允许做简易遮挡/置零实验；
    \item \textbf{V2b 后上。} 只有当 probe 路径允许对 branch 级变量反传并稳定记录梯度统计时，才启用 \eqref{eq:sensitivity}--\eqref{eq:sbranch_v2b}；
    \item 若连 V2a 的 activation/output proxy 都暂不可用，则直接把 \(w_b\) 置零，保持算法退化为 V1/V3。
\end{enumerate}

\paragraph{为什么这是“软约束”而不是“硬均衡”}
我们并不要求 \(A_o=A_m=A_\Delta\)。不同 slice 对三路的需求本就不同。我们只惩罚“单路塌缩”，不追求“完全平均”。因此 \(S_{branch}\) 只能作为辅助项，不能盖过 hardest-slice 主目标。

\subsection{成本感知总评分}
将以上各项合并，总评分写为
\begin{equation}
\begin{aligned}
S_{raw}(g)=L(g)\Big[&w_q Q_{hard}(g)+w_h R_{hard}(g)+w_c C_{sel}(g)+w_e P_{easy}(g) \\
&+w_i M_{ins}(g)+w_p LP(g)+w_b S_{branch}(g)+w_n N(g)\Big]-w_d D(g),
\end{aligned}
\label{eq:sraw}
\end{equation}
其中：
\begin{itemize}[leftmargin=2.2em,itemsep=0.15em]
    \item \(N(g)=1/\sqrt{n_{seen}(g)+1}\) 表示新颖度；
    \item \(D(g)\) 表示文件生成失败、路径错误、空轨迹等工程惩罚；
    \item \(w_q,w_h,w_c,w_e,w_i,w_p,w_b,w_n,w_d\) 为可调权重；
    \item 若 branch-aware 或 \(p(a=1)\) 统计暂不可用，则应显式把对应项关闭，而不是保留公式却在实现中隐式跳过。
\end{itemize}
这里推荐满足
\[
w_q \ge w_h \gg w_i,
\]
也就是：\textbf{hardest slice 总体正确性至少不应比 action1 恢复更轻，而 insertion 只能作为弱辅助项或纯监控项。} 这样才能避免总分被“多做 action1”单独主导。

考虑训练成本后，真正用于 outer 选择的效用定义为
\begin{equation}
U(g)=\frac{S_{raw}(g)}{Cost(g)+\varepsilon_u}.
\label{eq:utility}
\end{equation}
这一步非常关键，因为你们的外层瓶颈不是“网络更新慢”，而是“每个 outer 样本都贵”。

\section{版本路线：从简单到复杂逐步实现}

\subsection{V0：PLR-Lite（最小可替代版）}
\subsubsection*{目标}
用最小风险替代当前高耗时 outer RL。V0 不加 branch-aware，也不加两阶段预算；只保留以下三件事：
\begin{enumerate}[leftmargin=2.2em,itemsep=0.25em]
    \item learnability gate \eqref{eq:learnability}；
    \item learning progress \eqref{eq:lp}；
    \item replay 更新 \eqref{eq:q_update} 与优先级采样 \eqref{eq:priority_sampling}。
\end{enumerate}

V0 的总分数定义为
\begin{equation}
S_{V0}(g)=L(g)\big[w_\Delta LP(g)+w_J J(g)+w_n N(g)\big]-w_d D(g).
\label{eq:sv0}
\end{equation}

\subsubsection*{适用场景}
\begin{itemize}[leftmargin=2.2em,itemsep=0.2em]
    \item 先替换 EDRL/RARL 这类更重 outer 控制器；
    \item 验证“轻量 score-based selector + replay”就能跑通；
    \item 作为后续所有版本的保底退化路径。
\end{itemize}

\subsubsection*{推荐参数}
\begin{center}
\begin{tabular}{p{4.0cm}p{3.0cm}p{6.4cm}}
\toprule
参数 & 推荐值 & 说明 \\
\midrule
\(c\) & 0.55 & learnability 中心，落在“过易与过难之间”的中间偏低位置 \\
\(\sigma_L\) & 0.18 & 可学习带宽 \\
\(\alpha_q\) & 0.9 & priority exponent \\
\(\beta_q\) & 0.20 & score EMA \\
\(\varepsilon_q\) & 1e-4 & 防止零概率 \\
\(w_\Delta,w_J,w_n,w_d\) & 0.45, 0.25, 0.15, 0.15 & 先偏向 learning\ progress \\
replay/new & 0.65 / 0.35 & 初始混合 \\
buffer size & 64 & 小而稳，便于快速周转 \\
\bottomrule
\end{tabular}
\end{center}

\subsection{V1：SR-PLR（Severity-aware Removal PLR）}
\subsubsection*{目标}
从“通用 high-value replay”升级为“定向 hardest slice replay”。V1 的关键不是简单追求更多 \(a=1\)，而是把 score 从全局统计改成针对 hardest removal slice 的\textbf{边界决策质量}统计。

\subsubsection*{总分数}
\begin{equation}
S_{V1}(g)=L(g)\big[w_q Q_{hard}(g)+w_h R_{hard}(g)+w_c C_{sel}(g)+w_e P_{easy}(g)+w_i M_{ins}(g)+w_p LP(g)+w_n N(g)\big]-w_d D(g).
\label{eq:sv1}
\end{equation}

\subsubsection*{核心变化}
\begin{enumerate}[leftmargin=2.2em,itemsep=0.25em]
    \item \textbf{用 \(Q_{hard}+R_{hard}\) 取代通用 difficulty。} 优先修复 \(sev\ge 5\) 的 removal 边界，而不是全局最难分布；
    \item \textbf{引入 \(C_{sel}\)，但允许回退。} 若 \(p(a=1)\) 暂不可稳定导出，则使用 \(\tilde C_{sel}\) 或直接关闭该项；
    \item \textbf{引入 \(P_{easy}\)。} 保住 easy slice 的校准与稳定性，避免为了 hardest slice 而整体变得激进；
    \item \textbf{insertion 只做弱辅助。} \(M_{ins}\) 只占很小权重，第一版甚至可以设为 0，仅保留监控。
\end{enumerate}

\subsubsection*{推荐参数}
\begin{center}
\begin{tabular}{p{4.0cm}p{3.0cm}p{6.4cm}}
\toprule
参数 & 推荐值 & 说明 \\
\midrule
\(\tau_H\) & 5 & hard severity threshold \\
\(\tau_E\) & 3 & easy threshold \\
\(\lambda_{fp}\) & 0.20 & easy slice 假阳性惩罚 \\
\(w_q,w_h,w_c,w_e,w_i,w_p,w_n,w_d\) & 0.28, 0.18, 0.14, 0.15, 0.05, 0.10, 0.05, 0.05 & 让 hardest 总体正确性重于 action1 恢复，insertion 仅弱辅助 \\
\(w_q:w_h:w_i\) & \(0.28:0.18:0.05\) & 满足 \(w_q \ge w_h \gg w_i\)，避免主目标被“多做 action1”带偏 \\
\bottomrule
\end{tabular}
\end{center}

\subsection{V2：SBR-PLR（Stream- and Branch-aware PLR）}
\subsubsection*{目标}
V2 解决“outer 调度与 v3 主干脱节”的问题，但前提是 branch telemetry 至少部分可见。除保留 V1 所有项外，新增 \(S_{branch}(g)\) 作为辅助项。这里进一步分成两个现实层级：
\begin{enumerate}[leftmargin=2.2em,itemsep=0.2em]
    \item \textbf{V2a：proxy-only branch-aware。} 只依赖 activation-level / output-level proxy，是当前更现实的第一实现版本；
    \item \textbf{V2b：gradient-aware branch-aware。} 只有在 telemetry 打通后才启用，是后续增强版。
\end{enumerate}
若连 V2a 的 proxy 都暂不可稳定导出，则 V2 应继续关闭，不应强行上线。

\subsubsection*{总分数}
\begin{equation}
S_{V2}(g)=S_{V1}(g)+L(g)\,w_b S_{branch}(g).
\label{eq:sv2}
\end{equation}

\subsubsection*{关键思想}
\begin{itemize}[leftmargin=2.2em,itemsep=0.25em]
    \item 如果某个 level 只能把 \(z^{(o)}\) 打热，而 \(meta\) 与 \(\Delta o\) 路分支几乎无活性或无敏感度，则它更像“单路刷分 level”，不应被过度优先；
    \item 如果某个 level 在 hardest slice 上使三路都产生非零可学习信号，说明它更能训练出稳健边界；
    \item 这样外层选择的不只是“短期提分 level”，而是“更能训练 v3 结构先验的 level”；
    \item 但这条思路只有在 branch proxy 与真实性能确实相关时才成立，因此 V2 必须默认带退化开关；
    \item 工程上应先上 V2a，再决定是否值得升级到 V2b。
\end{itemize}

\subsubsection*{推荐参数}
\begin{center}
\begin{tabular}{p{4.0cm}p{3.0cm}p{6.4cm}}
\toprule
参数 & 推荐值 & 说明 \\
\midrule
\(w_b\) & 0.05 -- 0.10 & 只做辅助项；V2a 建议从 0.05 起，V2b 再视效果上调 \\
\(\eta_b\) & 0.35 & 分支失衡惩罚 \\
\(\varepsilon_a,\varepsilon_\delta,\varepsilon_G,\varepsilon_b\) & 1e-6 & 数值稳定 \\
V2a 统计样本 & hardest slice only & 只在 \(sev\ge 5\), removal 上统计 branch proxy \\
V2b 梯度采样频率 & 每个 probe 末 1 次 & 降低额外开销；若无法稳定反传，则停留在 V2a \\
\bottomrule
\end{tabular}
\end{center}

\paragraph{保守开关机制}
为避免 \(S_{branch}\) 误导总分数，推荐加入保守门控：若最近 \(m\) 轮验证中，\(S_{branch}\) 与 \(\Delta J\) 的秩相关系数不为正，则暂时将 \(w_b\) 退火到 0。即
\begin{equation}
 w_b(t+1)=
 \begin{cases}
 \rho_w\, w_b(t), & \text{if } \mathrm{Corr}(S_{branch},\Delta J) \le 0,\\
 \min(w_b^{max}, w_b(t)/\rho_w), & \text{otherwise},
 \end{cases}
\label{eq:wb_schedule}
\end{equation}
其中 \(\rho_w=0.5\) 推荐。这样即便 branch-aware 信号暂时无效，也能自动退化回 V1。

\subsection{V3：BSR-PLR（Budgeted Severity-aware Replay PLR）}
\subsubsection*{目标}
V3 是\textbf{最有希望明显降时}的关键版本。它加入两阶段预算控制，把大部分无效 full evaluation 换成廉价 probe。这里的“更快”首先指外层 candidate-evaluation 成本更低，而不是先验承诺整个实验 wall-clock 必然按同一比例下降。

\subsubsection*{两阶段流程}
\paragraph{阶段 A：cheap probe}
对每个候选 level 只做少量 rollout 或极短 inner short-train，得到
\[
J_{frozen}(g),\; LP(g),\; Q_{hard}(g),\; R_{hard}(g),\; C_{sel}(g),\; P_{easy}(g),\; M_{ins}(g),\; S_{branch}(g),\; D(g).
\]
基于这些量计算 \(U(g)\)。

\paragraph{阶段 B：full short-train}
只对满足 \(L(g)\ge \tau_L\) 的候选中、\(U(g)\) 前 \(M\) 个 level 做正式短训，并用训练后的真实结果更新 \(q_g\)。

\subsubsection*{预算函数}
为避免把大预算浪费在“太差但其实不可学”的候选上，先定义 learnability 硬门
\begin{equation}
\Gamma_L(g)=\I\big(L(g)\ge \tau_L\big).
\label{eq:learnability_gate_budget}
\end{equation}
记 full short-train 预算为 \(B(g)\)，定义
\begin{equation}
B(g)=
\begin{cases}
0, & L(g)<\tau_L,\\
B_{min}+(B_{max}-B_{min})\cdot \sigmoid\Big(\alpha_B(\tau_U-U(g)) + \beta_B(\tau_R-R_{hard}(g))\Big), & L(g)\ge\tau_L.
\end{cases}
\label{eq:budget}
\end{equation}
解释如下：
\begin{itemize}[leftmargin=2.2em,itemsep=0.2em]
    \item 若 \(L(g)<\tau_L\)，则该候选\textbf{直接跳过 full evaluation}，因为它已经落到“太难或不可学”的一侧；
    \item 若 \(U(g)\) 高、\(R_{hard}(g)\) 也高，说明该 level 已经学得不错或 probe 已充分，预算靠近 \(B_{min}\)；
    \item 若 \(U(g)\) 尚可但 \(R_{hard}(g)\) 仍差，说明这是“仍值得进一步修 hardest slice”的 level，预算靠近 \(B_{max}\)。
\end{itemize}

\subsubsection*{成本感知效用}
V3 采用 \eqref{eq:utility}，其中
\begin{equation}
Cost(g)=c_0 + c_1\,T_{probe}(g)+c_2\,B(g),
\label{eq:cost}
\end{equation}
\(T_{probe}(g)\) 为 probe 时间或统一设定的 probe 步数。

\subsubsection*{推荐参数}
\begin{center}
\begin{tabular}{p{4.0cm}p{3.0cm}p{6.4cm}}
\toprule
参数 & 推荐值 & 说明 \\
\midrule
候选数 \(N_{cand}\) & 6 & 每轮只看 6 个候选，控制外层成本 \\
full 进入数 \(M\) & 2 & 只对前 2 个做正式短训 \\
\(B_{min},B_{max}\) & 0.25, 1.00 & 相对当前 full short-train 标准预算的比例 \\
probe 预算 & 0.12 -- 0.18 倍 full & 一般推荐 0.15 倍 \\
\(\tau_L\) & 0.30 -- 0.40 & learnability 硬门，建议先从 0.35 起步 \\
\(\alpha_B,\beta_B\) & 3.0, 4.0 & 让 hardest-recovery 主导预算分配 \\
\(\tau_U,\tau_R\) & 0.55, 0.20 & utility 与 hard-recovery 的中间门槛 \\
\bottomrule
\end{tabular}
\end{center}

\paragraph{理论上的加速来源}
若当前 PLR\_UED 每轮对 \(N_{cand}\) 个候选都做 full evaluation，则其每轮成本近似为
\[
T_{old}\approx N_{cand}\,C_{full}.
\]
V3 后变为
\[
T_{new}\approx N_{cand}\,C_{probe} + M\,C_{full},
\]
若取 \(N_{cand}=6\), \(M=2\), \(C_{probe}=0.15C_{full}\)，则
\[
\frac{T_{new}}{T_{old}}\approx \frac{6\times0.15+2}{6}=0.4833.
\]
这说明：若只看\textbf{候选评估型训练开销}，V3 在结构上有机会把该部分压到旧版本的约 \(48\%\)。但这\textbf{不等于}最终 wall-clock 必然减半，因为真实总耗时还包含事件文件生成、子进程启动、日志落盘、checkpoint 切换等额外成本。因此 V3 的速度收益必须配合 runtime profiling 报告，不能只靠公式口径下结论。

\subsection{V4：SABER-PLR（完整版）}
\subsubsection*{目标}
V4 在 V3 的基础上，把“找对 level”这件事做得更稳。它新增三项机制：
\begin{enumerate}[leftmargin=2.2em,itemsep=0.25em]
    \item 邻域编辑课程：从 top replay level 出发做局部扰动，不再只重放单点；
    \item adaptive new/replay mixing：根据当前 hardest slice 修复情况调节探索比例；
    \item 轻量 contextual bandit/UCB 选择器：优先复用或轻改现有 UCB/TS 式候选选择器，而不是再引入一套重 outer RL。
\end{enumerate}

\subsubsection*{邻域编辑}
从 replay top-\(K\) level 中抽一个种子 \(g_i\)，做局部编辑：
\begin{equation}
 g' = (\mu_A+\delta_A,\; \mu_B+\delta_B,\; p+\delta_p,\; n+\delta_n,\; pattern'),
\label{eq:edit}
\end{equation}
其中推荐默认扰动集合：
\[
\delta_A,\delta_B\in\{\pm 10\},\qquad \delta_p\in\{\pm 0.1\},\qquad \delta_n\in\mathcal{N}^{adj}_{files}(n).
\]
其中 \(\mathcal{N}^{adj}_{files}(n)\) 表示\textbf{当前 generator 所支持的离散相邻文件规模档位}，而不是默认存在连续的 \(\pm 100\) 扰动。对 \(pattern\) 也同理：第一版只能在\textbf{当前 generator 真正支持}的模式集合内切换；若当前实现只稳定支持少数模式，则不得把 \(ab\leftrightarrow aba\leftrightarrow abba\) 这类编辑写成现成能力。

\subsubsection*{adaptive new/replay mixing}
定义最近窗口内 hardest-slice 修复率为
\begin{equation}
\rho_{hard}(t)=\frac{1}{W}\sum_{k=t-W+1}^{t} \E\big[\I(sev\ge 5, str=\texttt{removal}, a=1)\,r\big].
\label{eq:rhohard}
\end{equation}
再定义当前新样本概率：
\begin{equation}
p_{new}(t)=\clip\Big(p_0 + k_1(H^*-\bar H_t)+k_2(\rho_{hard}^*-\rho_{hard}(t)),\; p_{min},\; p_{max}\Big),
\label{eq:pnew}
\end{equation}
其中 \(\bar H_t\) 可为 recent entropy 或 recent diversity 指标。含义是：
\begin{itemize}[leftmargin=2.2em,itemsep=0.2em]
    \item 若动作熵过低、hardest recovery 仍弱，则提高 \(p_{new}\) 去找新 level；
    \item 若 hardest recovery 已起色但不稳定，则提高 replay 比例做巩固训练。
\end{itemize}

\subsubsection*{轻量选择器：UCB over utility}
不训练重 outer RL，只在当前候选集上做轻量 UCB：
\begin{equation}
A(g)=U(g)+\beta_{ucb}\sqrt{\frac{\log(t+1)}{n(g)+1}},
\label{eq:ucb}
\end{equation}
每轮取 \(\argmax_g A(g)\) 作为进入 full short-train 的主候选之一。

\subsubsection*{推荐参数}
\begin{center}
\begin{tabular}{p{4.0cm}p{3.0cm}p{6.4cm}}
\toprule
参数 & 推荐值 & 说明 \\
\midrule
replay top-\(K\) & 8 & 从 top-8 中采样邻域种子 \\
邻域编辑比例 & 0.5 & 新候选中一半来自邻域编辑 \\
\(p_0,p_{min},p_{max}\) & 0.35, 0.15, 0.65 & 新样本基础概率 \\
\(k_1,k_2\) & 0.20, 0.25 & 熵与 hardest-recovery 双调节 \\
\(\beta_{ucb}\) & 0.05 -- 0.12 & 轻量探索系数 \\
窗口 \(W\) & 10 & 统计 \(\rho_{hard}(t)\) 的滑窗长度 \\
\bottomrule
\end{tabular}
\end{center}

\section{统一训练阶段设计}

\subsection{Stage 0：inner backbone 预热}
保持你们现有的 Phase1，不做任何变动：
\begin{enumerate}[leftmargin=2.2em,itemsep=0.25em]
    \item 用当前 base 分布训练出稳定的 \(\theta_{base}\)；
    \item 缓存基础统计：整体 reward、action ratio、trace 切片、hardest slice 基线；
    \item 初始化 replay buffer，可以把 base 附近的一小部分分布作为“热启动候选”。
\end{enumerate}

\subsection{Stage 1：候选生成}
每一轮 outer 先构造候选集 \(\mathcal{C}_t\)，由三部分组成：
\begin{equation}
\mathcal{C}_t = \mathcal{C}^{replay}_t \cup \mathcal{C}^{edit}_t \cup \mathcal{C}^{new}_t.
\label{eq:candidates}
\end{equation}
其中：
\begin{itemize}[leftmargin=2.2em,itemsep=0.2em]
    \item \(\mathcal{C}^{replay}_t\)：从优先级 buffer 中抽的历史高价值 level；
    \item \(\mathcal{C}^{edit}_t\)：从 replay top level 生成的邻域编辑 level；
    \item \(\mathcal{C}^{new}_t\)：按当前 \(p_{new}(t)\) 抽的全新 level。
\end{itemize}

\subsection{Stage 2：cheap probe 与评分}
\paragraph{实现前提}
在进入这一阶段之前，至少要保证 probe/trace 链路能稳定导出 \(severity\)、stream、语义动作、reward；若 \(p(a=1)\) 或 branch 统计暂不可用，则对应项必须显式关闭，而不是在文档里保留、在实现里默默跳过。

对 \(\mathcal{C}_t\) 中每个候选 level：
\begin{enumerate}[leftmargin=2.2em,itemsep=0.25em]
    \item 生成对应事件文件；
    \item 用冻结或当前 backbone 做极短 rollout / short-train；
    \item 只在 trace 中抽取 hardest 与 easy slice；
    \item 计算 \(L, LP, Q_{hard}, R_{hard}, C_{sel}, P_{easy}, M_{ins}, S_{branch}, D\)；
    \item 算出 \(S_{raw}(g)\) 与 \(U(g)\)。
\end{enumerate}

\subsection{Stage 3：选择与 full short-train}
\begin{enumerate}[leftmargin=2.2em,itemsep=0.25em]
    \item 先按 \(\Gamma_L(g)=\I(L(g)\ge \tau_L)\) 过滤掉不可学候选，再用 \eqref{eq:ucb} 或直接按 \(U(g)\) 排序，选出前 \(M\) 个候选；
    \item 对每个被选中的候选，根据 \eqref{eq:budget} 分配预算 \(B(g)\)；
    \item 在该候选上做正式 short-train，得到真实 \(J(g)\) 与 trace 统计；
    \item 用真实结果更新 \(q_g\)、\(n_{seen}(g)\)、\(\rho_{hard}(t)\)、\(p_{new}(t)\)。
\end{enumerate}

\subsection{Stage 4：回放更新与切换条件}
若最近 \(W\) 轮满足以下任一条件，可进入下一阶段或触发终止：
\begin{enumerate}[leftmargin=2.2em,itemsep=0.25em]
    \item \(\rho_{hard}(t)\) 持续上升且方差收缩；
    \item \(R_{hard}(g)\) 进入平台区，但 easy slice 未明显受损；
    \item replay top level 已出现稳定的邻域簇，而非单点垄断；
    \item 达到外层预算上限。
\end{enumerate}

\section{推荐默认参数与调参顺序}

\subsection{推荐默认参数总表}
\begin{center}
\small
\begin{longtable}{p{3.2cm}p{2.8cm}p{7.2cm}}
\toprule
类别 & 参数 & 推荐值 \\
\midrule
\endfirsthead
\toprule
类别 & 参数 & 推荐值 \\
\midrule
\endhead
learnability & \(c,\sigma_L\) & \(0.55,\;0.18\) \\
priority replay & \(\alpha_q,\beta_q,\varepsilon_q\) & \(0.9,\;0.20,\;10^{-4}\) \\
hard/easy slice & \(\tau_H,\tau_E\) & \(5,\;3\) \\
selective action & \(\lambda_{fp}\) & \(0.20\) \\
stream asymmetry & \(w_i\) / monitor & \(w_i=0.05\) 或直接取 0，仅把 insertion 作为弱辅助或监控项 \\
score weights (V1) & \makecell[l]{$w_q,w_h,w_c,w_e,$\\$w_i,w_p,w_n,w_d$} & \makecell[l]{0.28, 0.18, 0.14, 0.15\\0.05, 0.10, 0.05, 0.05} \\
branch-aware & \(w_b,\eta_b\) & \(0.05\text{--}0.10,\;0.35\)；V2a 先上，V2b 后上 \\
probe/full & \(N_{cand},M\) & \(6,\;2\) \\
budget & \(B_{min},B_{max},\tau_L\) & \(0.25,\;1.00,\;0.35\) \\
probe ratio & \(C_{probe}/C_{full}\) & \(0.15\) \\
budget slope & \(\alpha_B,\beta_B\) & \(3.0,\;4.0\) \\
new/replay & \(p_0,p_{min},p_{max}\) & \(0.35,\;0.15,\;0.65\) \\
adaptive mix & \(k_1,k_2\) & \(0.20,\;0.25\) \\
UCB & \(\beta_{ucb}\) & \(0.08\) \\
edit neighborhood & \(\delta_\mu,\delta_p,\delta_n\) & \(\pm10,\;\pm0.1,\;\mathcal{N}^{adj}_{files}(n)\) \\
window & \(W\) & 10 \\
\bottomrule
\end{longtable}
\end{center}

\subsection{建议的实现顺序}
为降低工程风险，建议按下面顺序推进，而不是一次性上完整版：
\begin{enumerate}[leftmargin=2.2em,itemsep=0.3em]
    \item \textbf{先做 V0。} 目标是用最小替换跑通 replay + score-selector 的主链路；
    \item \textbf{再做 V1。} 一旦 trace 级 hardest slice 指标能稳定导出，就立刻切到定向 score；
    \item \textbf{再做 V3。} 因为你当前最痛的是时间，预算控制应早于 branch-aware 上线；
    \item \textbf{然后再决定是否做 V2。} 若 branch proxy 已可稳定导出，先做 V2a；只有当 V2a 确实有效且 telemetry 打通后，再考虑 V2b；
    \item \textbf{最后做 V4。} 邻域编辑与 adaptive mixing 放到稳定主链之后，不要一开始全上。
\end{enumerate}

如果只能做三版，我建议依次做：
\[
\text{V0} \rightarrow \text{V1} \rightarrow \text{V3}.
\]
因为这三版已经覆盖：\textbf{替换重 outer、定点修 hardest slice、显著降时} 这三件核心事情。

\section{为什么 branch-aware 项值得加，但不能当主角}
\subsection{值得加的理由}
若没有 \(S_{branch}(g)\)，新算法仍然可能有效，但它更像“PLR\_UED 外层 + v3 内层”的并列组合。加入 \(S_{branch}(g)\) 后，outer 在挑分布时会回答一个更细的问题：
\begin{quote}
\textbf{哪些 level 不仅能提分，而且在 hardest slice 上更可能训练出三路可学习而非单路塌缩的表征？}
\end{quote}
这样就把 v3 的结构先验抬升成了 outer 的候选选择依据。但必须强调：这是一种\textbf{研究扩展价值}，不是当前项目里已经稳定上线的能力。工程上更合理的做法是先验证 V2a 的 proxy-only 版本，再决定是否值得走到 V2b。

\subsection{不能当主角的理由}
当前已有证据说明，hardest OOD 的主导因素仍是 \textbf{训练分布调度}，而不是继续加重 inner 表征。因此：
\begin{itemize}[leftmargin=2.2em,itemsep=0.2em]
    \item \(S_{branch}\) 应作为辅助项，权重不宜过大；
    \item 一旦发现 \(S_{branch}\) 与真实提升脱钩，就应通过 \eqref{eq:wb_schedule} 自动退化；
    \item 即使最终发现 branch-aware 项收益有限，V2 也应稳定退化回 V1，V3 仍可独立作为速度主线存在。 
\end{itemize}

\section{实验建议、消融建议与预期结论口径}

\subsection{最关键的对比与消融}
推荐最少做以下对比：
\begin{enumerate}[leftmargin=2.2em,itemsep=0.25em]
    \item PPO vs PPO\_NEW-v3/v3.1 vs PLR\_UED vs V1 vs V3；
    \item 在 \(O\_{10,90}\) 上至少做 3 seeds，并强制报告 trace-level slice，而不是只报平均回报；
    \item V1 去掉 \(Q_{hard}\)（验证 hardest-slice 总体正确性是否必要）；
    \item V1 去掉 \(C_{sel}\)（验证 selective action1 是否必要）；
    \item V1 令 \(w_i=0\)（验证 insertion 弱辅助是否真的有价值）；
    \item V1 去掉 \(P_{easy}\)（验证 easy-slice 保护是否必要）；
    \item V2a vs V2b（验证 proxy-only 版本是否已经足够）；
    \item V2 去掉 \(S_{branch}\)（验证 branch-aware 是否真的增益）；
    \item V3 去掉 learnability gate 或去掉两阶段预算（验证主要加速来源）；
    \item V4 去掉邻域编辑（验证是否避免只记单点）。
\end{enumerate}

\subsection{推荐主指标}
不建议只报 average reward。至少要同时报告：
\begin{itemize}[leftmargin=2.2em,itemsep=0.2em]
    \item hardest slice overall quality：\(Q_{hard}\)；
    \item hardest slice recovery：\(R_{hard}\)；
    \item selective action quality：\(C_{sel}\)；
    \item easy-slice preservation：\(P_{easy}\)；
    \item insertion weak monitor：\(M_{ins}\)；
    \item 总体平均回报；
    \item 外层 wall-clock、单轮候选评估耗时以及 \(\Gamma_L(g)=0\) 的跳过比例；
    \item replay 覆盖率、top-k 邻域覆盖率。
\end{itemize}

\subsection{应当如何写“理论安全性”}
不能诚实地声称“真实测试回报必定提升”，因为这需要超出当前信息的严格证明。但可以较稳地表述为：
\begin{enumerate}[leftmargin=2.2em,itemsep=0.25em]
    \item 新算法是对现有 PLR 主线的\textbf{保守增强}，而非推翻；
    \item 若新增 hardest-slice 或 branch-aware 信号无效，权重可自动退化到 0，算法退化为 V0/V1；
    \item 两阶段预算控制对外层\textbf{候选评估型训练开销}有明确的削减机制，因此“更快”可以给出结构性解释；但这不是对最终 wall-clock 减半的直接保证；
    \item hardest-slice 定向 score 是对当前 trace 证据的直接利用，因此比通用外层难度目标更贴近真实失效模式。
\end{enumerate}

\section{总结：建议你实际落地时如何推进}
\begin{tcolorbox}[colback=blue!3,colframe=MidnightBlue,title=一句话总结]
\textbf{不要再把主要创新放在更复杂的 outer RL 控制器上，而要把现有 PLR/UED 提炼成一个“定点修复 severity 5/6 removal 边界、以 V1/V3 为主线、先做可验证的速度优化、branch-aware 仅作后续可选增强”的高样本效率课程重放框架。}
\end{tcolorbox}

更具体地说：
\begin{enumerate}[leftmargin=2.2em,itemsep=0.3em]
    \item 第一批实现建议先做 V0，快速替掉重 outer RL；
    \item 第二批直接切到 V1，把 score 对齐到 hardest slice；
    \item 若时间压力很大，第三批优先做 V3，而不是先做 V2；
    \item branch-aware 项值得做，但应视为“在 telemetry 打通后再验证”的辅助增强，而不是第一优先级；
    \item 完整版 V4 适合作为论文中的主方法，但工程实现上应在 V1/V3 验证稳定后再上。
\end{enumerate}

\vspace{0.8em}
\begin{tcolorbox}[colback=green!3,colframe=OliveGreen,title=推荐命名]
\textbf{工程主线名：SR-PLR / BSR-PLR}\\
中文：\textbf{严重度感知的高效优先级课程重放}\\
若后续 branch-aware 扩展被真实验证，再把完整版命名为：\textbf{SABER-PLR}。这样命名更符合“先落地 V1/V3，再决定是否引入 V2/V4”的实现顺序。
\end{tcolorbox}

\end{document}
